\section{Parent adaptation and match evaluation}\label{app:parent-adaptation}
\subsection{Policy and value functions}
We initialize the three phase networks from the original RL-WP weights.
Their inherited raw action scores $q_\theta(o,a)$ are not probabilities.
For the legal action set $\mathcal A(o)$, training samples from
\begin{equation}
 \pi_{\theta,T}(a\mid o)=
 \frac{\exp(q_\theta(o,a)/T)}
 {\sum_{b\in\mathcal A(o)}\exp(q_\theta(o,b)/T)},\qquad a\in\mathcal A(o).
\end{equation}
The public endgame guard can further narrow $\mathcal A(o)$. At evaluation,
the agent chooses a maximizing legal score, so a positive temperature does not
change the mathematical argmax for fixed scores. Temperature can nevertheless
change the learned policy through sampling and optimization. Dropout remains
disabled during both collection and gradient updates, preserving a consistent
policy for stored log probabilities.

The inherited actors retain the upstream 256-wide, two-layer Transformer
architecture. Native features reproduce the original 300-by-50 representation,
including its historical input conventions. A separate value network receives
a 560-dimensional observation of its own hand and public state, using two
256-unit tanh hidden
layers followed by a scalar output. Its output layer is initialized to zero;
it shares no actor parameters. Both actor and critic use only information available to the acting player.

\subsection{Player-relative returns}
Both players use the same current stochastic policy during self-play. Rewards
are zero until match termination, when the acting player receives $+1$, $0$,
or $-1$ for a win, draw, or loss. We use an undiscounted, player-relative
adaptation of generalized advantage estimation \cite{schulman2015gae}.
Let $d_t$ indicate match termination and let $s_t$ be $+1$ when consecutive
decision states have the same acting player and $-1$ otherwise. Then
\begin{align}
 \delta_t &= r_t+(1-d_t)s_t V(o_{t+1})-V(o_t),\\
 A_t &= \delta_t+(1-d_t)s_t\lambda A_{t+1},\qquad \lambda=0.95.
\end{align}
Rollout truncation bootstraps from the next observation. Round boundaries
retain continuation value unless the match ends. A terminal transition masks
the value of the automatically reset next match. The sign follows the actual
next acting player, which need not alternate at every decision phase.

Advantages are normalized over each rollout batch before optimization.
We optimize the clipped PPO surrogate \cite{schulman2017ppo} with clipping
parameter 0.2, add
one-half the value mean-squared error, and subtract 0.02 times policy entropy.
The global trainable gradient norm is clipped at 0.5. Adam uses
$\epsilon=10^{-5}$ and separate actor/value learning-rate groups. The first
20 updates freeze actor parameters and fit only the critic; the next 500
updates train both. This warmup is part of the reported budget.

\subsection{Fixed development comparison}
Each update collects 128 decisions from each of 128 parallel environments and
uses three minibatch epochs with minibatch size 512. Each model therefore
receives exactly 8,519,680 simulator decisions in 520 updates. The value
learning rate is $10^{-4}$. The warm setting uses $T=0.05$ and actor learning
rate $3\times10^{-6}$; the cold setting uses $T=0.01$ and
$6\times10^{-7}$. Both settings are trained at seeds 6710001 and 6720001.
Only the final update-520 checkpoints enter the comparison. Intermediate
snapshots are retained for accounting and recovery, not strength selection.

The comparison measures the joint effect of temperature and actor learning rate,
with their ratio held fixed. The settings can produce different sampled
trajectories, entropy and adaptive optimizer states. Execution is eager float32
with TF32 disabled. The GPU is an
NVIDIA GeForce RTX 5070 Ti; the runtime records PyTorch 2.12.0+cu130 and NumPy
2.5.4. Evaluation is kept separate from training.

\subsection{Implementation and inference cost}
A shared C++/CUDA transition core runs the same rules on CPU and GPU. Differential
checks compare transitions and observation features with the Python reference and
published model inputs; checkpoints retain the complete state for exact restart.
The adapted actor retains greedy, search-free deployment. On a Ryzen 9 9900X
with two CPU threads, the canonical model averages 1.585 ms per cached-state
decision, including feature construction and public assistance. The original
RL-WP agent averages 1.974 ms on the same corpus. Appendix~\ref{app:systems}
reports the validation corpus, separate training budgets and timing procedure.

\subsection{Match-evaluation protocol}
\subsubsection{Duplicate deals and uncertainty}
For each seed we play two matches with candidate/opponent seats exchanged,
retaining the same seat-assigned deal streams and initial dealer. If the
candidate's two match scores are $w_{i0},w_{i1}\in\{0,1/2,1\}$, the observation
for inference is $x_i=(w_{i0}+w_{i1})/2$. The reported draw-adjusted win score
is $\bar x$. We resample whole pairs, not individual match legs, for 5,000
bootstrap replicates using random seed 99 and linearly interpolated 2.5th and
97.5th percentiles. Paired method comparisons resample the per-deal differences.
The intervals quantify deal variation for each fixed model. Results for the
two training seeds are reported separately.

Point difference is final candidate points minus opponent points, averaged
across the two seats. Because total points remain 60, this is twice the
candidate's final-points-minus-initial-30 statistic used in the original paper.
Positive expected points are reported separately and cannot substitute for a
win-score criterion. Shared deals across models or opponents do not create
additional independent deal samples.

\subsubsection{Development and frozen final rules}
The development study uses 5,000 duplicate pairs with seeds
7,200,000--7,204,999 for the initial policy and four trained endpoints. An arm
qualifies only if both training seeds have an absolute win-score 95\% interval
lower bound above 0.5. A cold-over-warm benefit requires positive lower bounds
for both paired score differences. Selection chooses cold if it qualifies and
shows that replicated benefit; otherwise it chooses qualified warm, then
qualified cold, and otherwise no method. The canonical model always uses
training seed 6710001 of the selected arm. There is no best-seed ranking.

The selection rules were fixed before this study's outcomes. Earlier development
explored other methods and included an inconclusive two-seed warm screen at
1,000 pairs. The present comparison uses fresh training seeds and 5,000 pairs;
a separate final test provides confirmation after this development process.

The frozen final test evaluates both selected models against assisted RL-WP
and, separately, assisted RL-Point. Each of four jobs plays 10,000 duplicate
pairs on seeds 100,000,000--100,009,999: 80,000 matches on 10,000 unique deal
streams. The primary criterion requires both RL-WP interval lower bounds to
exceed 0.5. The secondary RL-Point criterion is recorded separately. There is
no change to the selected checkpoints or comparison budget during this final test.
All four jobs run to completion. Four two-thread CPU workers run concurrently;
a separate experiment measures isolated inference latency.

\subsection{Development match results}
The initialized policy reproduces the comparator's paired outcomes on all
5,000 development deals: a win score of 0.5 and a point difference of zero.
The trained endpoints then separate the two adaptation settings.

\begin{table}[ht]
\centering\small
\begin{tabular}{rlrrr}
\toprule
Training seed & Setting & Score (\%) & 95\% interval (\%) & Point diff. \\
\midrule
6710001 & warm & 51.805 & [51.090, 52.525] & 2.3802 \\
6710001 & cold & 50.755 & [50.135, 51.395] & 1.3046 \\
6720001 & warm & 51.500 & [50.765, 52.245] & 2.1546 \\
6720001 & cold & 50.565 & [49.930, 51.215] & 1.1602 \\
\bottomrule
\end{tabular}

\caption{Development results against assisted RL-WP. Every row uses the same
5,000 duplicate deal seeds. Point differences use candidate-minus-opponent units.}
\label{tab:development}
\end{table}

Both warm seeds qualify under the fixed development rule
(Table~\ref{tab:development}). Cold qualifies only at the first seed, and its
second interval includes 50\%. The paired differences favor warm on both
seeds (Table~\ref{tab:paired}); cold does not meet the replicated-benefit
criterion. The fixed rule therefore selects warm and the first-seed canonical
checkpoint. The bundled intervention changes exploration and optimizer step size together.

\begin{table}[ht]
\centering\small
\begin{tabular}{rrrr}
\toprule
Training seed & Cold minus warm (pp) & 95\% interval (pp) & Point diff. change \\
\midrule
6710001 & -1.050 & [-1.875, -0.200] & -1.0756 \\
6720001 & -0.935 & [-1.810, -0.075] & -0.9944 \\
\bottomrule
\end{tabular}

\caption{Paired development differences. Negative values favor warm.}
\label{tab:paired}
\end{table}

The earlier warm screen, using training seeds 6681001 and 6691001, produced
51.275\% [49.725, 52.850] and 51.075\% [49.475, 52.725] on 1,000 shared
development pairs. Neither cleared the original replicated criterion.
This earlier screen remains a separate exploratory result.

\subsection{Untouched final evaluation}
All four fixed evaluations completed, totaling 80,000 matches on the same
10,000 unique deal seeds. The primary RL-WP criterion passes; the replicated RL-Point criterion fails
(Table~\ref{tab:final}). Both models and all evaluation deals were fixed before
this test.

\begin{table}[ht]
\centering\small
\begin{tabular}{rlrrr}
\toprule
Training seed & Opponent & Score (\%) & 95\% interval (\%) & Point diff. \\
\midrule
6710001 & WP & 51.4800 & [50.970, 51.975] & 2.4203 \\
6720001 & WP & 51.4450 & [50.945, 51.953] & 2.0221 \\
6710001 & Point & 50.4925 & [49.955, 51.040] & -1.8918 \\
6720001 & Point & 51.2875 & [50.720, 51.852] & -1.9603 \\
\bottomrule
\end{tabular}

\caption{Frozen final results. Each row contains 10,000 duplicate pairs;
all rows share deal seeds. WP and Point denote assisted RL-WP and RL-Point.}
\label{tab:final}
\end{table}

Against RL-WP, both selected models' win-score intervals lie above 50\%.
Their mean point differences are also positive: 2.4203 [2.1224, 2.7133]
and 2.0221 [1.7159, 2.3171]. Against RL-Point, the canonical seed's score
interval includes 50\%, while the second seed's interval excludes it.
The method therefore fails the predeclared both-model secondary win criterion.
Moreover, both point-difference intervals are negative: -1.8918
[-2.2427, -1.5386] and -1.9603 [-2.3197, -1.6089]. This distinction between
match wins and expected points is part of the result. Seed 6710001 remains
the canonical model, regardless of the second seed's RL-Point result.

\subsection{Longer adaptation: a development comparison}
\label{sec:duration}
Two fresh original-WP learners test whether more optimization improves the
same architecture. Each receives 20 critic-only updates and 4,000 policy
updates, totaling 65,863,680 decisions. Its 500-policy-update checkpoint is an
exact prefix of that trajectory, including optimizer and environment state.
Learning rates, exploration, returns and greedy deployment are unchanged.
The two short, two long and two retained models play 98,304 matches on 8,192
shared development deals, with both seats and assisted WP/Point opponents.

Table~\ref{tab:duration-primary} compares the long endpoints with their short
prefixes and retained comparators. The primary endpoint weights WP and Point
equally. One simultaneous 95\% bootstrap family covers all four contrasts;
qualification requires every lower bound to exceed zero. Each resampled block
retains the two opponents' distinct deals, both seats and every policy.
Only the second seed's long-minus-short interval excludes zero, so the
replicated improvement criterion fails.

\begin{table}[htbp]
\centering\small
\begin{tabular}{lrr}
\toprule
Seed & Long minus short & Long minus retained \\
\midrule
1 & $+0.5798$ [$-0.3191$, $+1.4788$] & $+0.5249$ [$-0.3906$, $+1.4404$] \\
2 & $+0.8942$ [$+0.0005$, $+1.7878$] & $+0.6104$ [$-0.2847$, $+1.5054$] \\
\bottomrule
\end{tabular}

\caption{Longer-training development win changes in percentage points.
Intervals are simultaneous across these four comparisons; 10,000 bootstrap
replicates resample 4,096 compound blocks.}
\label{tab:duration-primary}
\end{table}

The additional training improves mixture point margins relative to both short
and retained models at both seeds: all four descriptive paired intervals are
positive. The gains over retained models are 0.7847 [0.3480, 1.2129] and
0.9548 [0.5186, 1.3823] points. Both long endpoints also have individual
win-score intervals above 50\% against each published opponent
(Table~\ref{tab:duration-absolute}). These secondary results describe the
fixed endpoints; they do not replace the primary win-improvement test.
Point margins against RL-Point remain negative on average. The canonical
player is retained and the reserved prospective confirmation is unused.
The experiment measures scaling with training effort, not a new algorithm.

\begin{table}[htbp]
\centering\small
\begin{tabular}{llrr}
\toprule
Seed & Opponent & Win score (\%) & Point margin \\
\midrule
1 & WP & $52.7283$ [$51.8982$, $53.5828$] & $+3.6338$ [$+3.1267$, $+4.1406$] \\
1 & Point & $50.9460$ [$50.0916$, $51.8127$] & $-1.1873$ [$-1.7510$, $-0.6294$] \\
2 & WP & $51.8066$ [$50.9644$, $52.6367$] & $+2.9602$ [$+2.4473$, $+3.4666$] \\
2 & Point & $52.2217$ [$51.3489$, $53.0945$] & $-0.3379$ [$-0.8889$, $+0.2339$] \\
\bottomrule
\end{tabular}

\caption{Long endpoints on development deals. Intervals are descriptive
marginal 95\%, not simultaneous across opponents or endpoints.}
\label{tab:duration-absolute}
\end{table}

\subsection{Point initialization: points and wins separate}
\label{sec:point-initialization}
Changing the initial policy offers another route to adaptation. At each of two
fresh seeds, we train three arms: original WP with divisor 0.05, original Point
with divisor 0.05, and original Point with calibrated phase-specific divisors.
Each receives the same 20 critic-only and 500 policy updates, or 8,519,680
decisions. Within a seed, initial critic, optimizer, environment and random
states match; the two Point arms also start from identical actor parameters.
Deployment uses the unchanged greedy policy and public endgame assistance.

Calibration matches Point's mean normalized action entropy to WP at 0.05 on
75 source groups from an existing random-play input corpus. The other 25 groups
are reserved for descriptive validation. Divisors are fixed before training;
no match outcomes enter calibration. This controls one property of the initial
sampling policy, while also changing policy-loss derivatives.

The six endpoints and two retained agents play 131,072 matches on 8,192 shared
development deals. The primary criterion requires calibrated Point to improve
equal-mixture win score over both WP-fixed and retained controls at both seeds.
All four simultaneous intervals include zero (Table~\ref{tab:point-primary}).
Calibrated-minus-fixed Point mixture wins change by $+0.7385$ and $-0.3387$
percentage points, so the temperature benefit does not replicate either.

\begin{table}[htbp]
\centering\small
\begin{tabular}{lrr}
\toprule
Seed & Calibrated minus WP fixed & Calibrated minus retained \\
\midrule
1 & $+0.2350$ [$-0.7334$, $+1.2034$] & $+0.4486$ [$-0.5250$, $+1.4222$] \\
2 & $+0.2716$ [$-0.6764$, $+1.2196$] & $-0.0366$ [$-0.9674$, $+0.8942$] \\
\bottomrule
\end{tabular}

\caption{Point-calibrated development win changes in percentage points.
One simultaneous 95\% family covers all four contrasts; 10,000 bootstrap
replicates resample 4,096 compound blocks preserving all policies and seats.}
\label{tab:point-primary}
\end{table}

Point initialization does improve point margins (Table~\ref{tab:point-points}).
Both Point arms exceed the WP-fixed and retained controls at both seeds,
against each opponent and their mixture: all 24 descriptive paired point-change
intervals have positive lower bounds. The effect therefore appears with both
fixed and calibrated sampling. Calibration itself reduces the second seed's
WP point margin relative to Point-fixed by 0.8909 [0.3691, 1.4089] points.

\begin{table}[htbp]
\centering\small
\begin{tabular}{llrr}
\toprule
Seed & Point initialization & Gain over WP fixed & Gain over retained \\
\midrule
1 & Fixed divisor & $+1.9663$ [$+1.5116$, $+2.4374$] & $+1.8054$ [$+1.3414$, $+2.2760$] \\
1 & Calibrated & $+1.9724$ [$+1.5005$, $+2.4347$] & $+1.8115$ [$+1.3519$, $+2.2717$] \\
2 & Fixed divisor & $+2.0848$ [$+1.6028$, $+2.5638$] & $+2.1075$ [$+1.6507$, $+2.5581$] \\
2 & Calibrated & $+1.7689$ [$+1.2904$, $+2.2479$] & $+1.7916$ [$+1.3292$, $+2.2450$] \\
\bottomrule
\end{tabular}

\caption{Equal-mixture point-margin improvements under Point initialization.
Intervals are descriptive marginal 95\%. Margins are candidate minus opponent;
they are twice the original paper's final-minus-initial point convention.}
\label{tab:point-points}
\end{table}

Against published RL-Point, the calibrated endpoints score 50.9155\%
[50.1219, 51.7212] and 50.8240\% [50.0305, 51.5930] in wins. Their mean point
margins are $-0.0498$ and $+0.0383$, with both intervals containing zero.
These descriptive matchup results do not establish the primary gain over
adapted controls, equivalence or win noninferiority. The canonical player is
retained and prospective confirmation remains unused.

A diagnostic fixed before outcome inspection tests the six endpoints on the
old validation inputs. At both seeds and in all four phases, calibrated Point
has higher entropy at its training divisor, but lower entropy at a common
0.05 divisor and larger raw score gaps than Point-fixed. Higher sampling
entropy persists on this corpus without a replicated greedy win benefit.
The diagnostic measures policy outputs on reused inputs, not causal mediation
or equality of visited states. Appendix~\ref{app:reproduction} records the
calibration, computation and complete result bindings.

\subsection{Outcome and margin auxiliary prediction}
\label{sec:score-auxiliary}
We compare three Point-initialized conditions at each of two fresh seeds:
ordinary PPO, auxiliary outcome prediction and auxiliary margin prediction.
The two prediction tasks share a 65-output head and initialization. Margin
supervision adds the conditional likelihood within the observed loss, draw or
win group. Each condition receives the same 20 warmup and 500 policy updates
and collects 12,800 auxiliary matches; the control performs no auxiliary
optimization. Appendix~\ref{app:score-auxiliary} gives the losses and procedure.

Margin supervision changes 10.87\% and 8.18\% of greedy choices relative to
outcome supervision on the fixed validation inputs. Those changes do not
establish a playing benefit. All six primary simultaneous win-change intervals
include zero (Table~\ref{tab:aux-primary}). All 24 descriptive win/point
contrasts against the no-auxiliary and outcome controls also include zero
across both seeds, both opponents and their mixture.

\begin{table}[htbp]
\centering\small
\begin{tabular}{llrr}
\toprule
Seed & Margin minus & Win change (pp) & Simultaneous 95\% interval \\
\midrule
1 & No auxiliary & $-0.1953$ & $[-1.1108, +0.7202]$ \\
1 & Outcome & $-0.0305$ & $[-0.9607, +0.8997]$ \\
1 & Retained WP parent & $+0.8209$ & $[-0.2281, +1.8699]$ \\
2 & No auxiliary & $+0.3510$ & $[-0.5326, +1.2345]$ \\
2 & Outcome & $-0.1709$ & $[-1.0401, +0.6983]$ \\
2 & Retained WP parent & $+0.6805$ & $[-0.3703, +1.7314]$ \\
\bottomrule
\end{tabular}

\caption{Equal-mixture development win changes from margin supervision. One simultaneous 95\%
family covers all six contrasts; 131,072 matches use shared deals and both
seats. The 10,000-replicate bootstrap resamples 4,096 compound blocks.}
\label{tab:aux-primary}
\end{table}

All three Point-initialized conditions exceed the retained WP parents in point
margin at both seeds and all opponent scopes: all 18 descriptive paired
intervals have positive lower bounds. Because the no-auxiliary control shares
this advantage, it does not identify an auxiliary-learning gain. The result
separates a point-margin advantage from the additional hypothesis
that finer terminal targets improve action selection at this training budget.
